\documentclass[10pt,a4paper]{amsart}
\usepackage[margin=19mm]{geometry}
\usepackage{amsmath,amssymb,mathtools}
\usepackage[scr=rsfs,cal=euler]{mathalfa}
\usepackage{xfrac}
\usepackage[unicode,hidelinks,bookmarks=false]{hyperref}
\newcommand{\ie}{i.e.\ }
\newcommand{\cf}{cf.\ }
\newcommand{\resp}{respectively\ }
\newcommand{\Subsection}{\subsection}
\providecommand{\nobreakdash}{\nobreak\hbox{-}}
\setlength{\emergencystretch}{2em}
\multlinegap0pt
\usepackage{mathtools}
\newtheorem{theorem}{Theorem}
\DeclareMathOperator{\vol}{vol}
\title{Free energy of the Coulomb gas in the determinantal case on Riemann surfaces}
\author{Lucas Bourgoin (IRMA)}
\date{Source: arXiv:2508.20598v5 (CC BY 4.0)}
\begin{document}
\maketitle

\noindent\emph{Source and licence.} Free energy of the Coulomb gas in the determinantal case on Riemann surfaces. Version arXiv:2508.20598v5 (CC BY 4.0), distributed by arXiv under the Creative Commons Attribution 4.0 International licence, http://creativecommons.org/licenses/by/4.0/, as recorded in the arXiv licence metadata for this exact version and shown on the abstract page https://arxiv.org/abs/2508.20598v5. Full licence notice: \texttt{licenses/arxiv-2508.20598-CC-BY-4.0.txt}.\par\smallskip

\noindent\textit{Editorial scope.} Counted unit: the article's theorem TheoremeModifPreFinal (source0.tex lines 642--650). The article's complete original proof of it is delivered with it (source0.tex lines 652--661, one \texttt{proof} environment of 10 lines). No mathematics is written, repaired or re-derived: the statement and the proof are the article's own lines, in the article's own notation, and no retained line is substituted or re-flowed.  Retained uncounted as the source-local context the proof needs: lines 528--530; lines 567--575; lines 598--603; lines 629--634.  Nothing was omitted from any retained span, so the package carries no omission record.  No retained line contains a citation, so the wrapper carries no bibliography and no bibliography entry.  No retained line contains a figure environment, an image-inclusion command or drawing code, so no image is delivered and the package declares no asset dependency.  Editorial formatting only, in addition: an AMS \texttt{amsart} preamble that restates the article's own declarations and macros used by the retained lines, namely the environments \texttt{theorem} on separate counters, together with the standard packages those lines need.  The retained environments print in this document as Theorem 1 in order of appearance, whereas the article prints the same items with the numbering of its own full document.  The complete native source of the article as distributed in the arXiv e-print for this exact version is bundled unchanged as \texttt{sources/184031/source0.tex}.
     \begin{equation}\label{lnBgk}
         \ln B_{g,k}=-k\ln2\pi+\frac{1-g}{4}c_{0}=-k\ln2\pi +\frac{1-g}{2} \left(\frac{1}{2}-12\zeta'(-1)-3\ln2\pi-\ln2 \right).
         \end{equation}

    \begin{theorem}\label{TheoremPartitionGénéral}
        With $L$ a line bundle over $M$ of degree $k=N+g-1$ and $h$ an admissible metric on $L$, the following exact equality holds: 
        \begin{equation}\label{EqExactGeneral}
            \ln\mathcal{Z}_{\rho_{Ar},{\rho_{can}},N}^{(\theta)}(V)=\ln \mathcal{Z}_{\rho_{Ar},{\rho_{can}},N}^{(\theta)}(0)+\ln \det\square_{L,\rho_{Ar},\hat{h}_{k,V}}-\ln \det\square_{L,\rho_{Ar},h}+k^{2}S_{2}(V,\rho_{Ar},h)+\frac{k}{2}S_{1}(0,V,\rho_{Ar},h)
        \end{equation}
        with
        \begin{equation}\label{EqExactGeneralV=0}\ln\mathcal{Z}_{\rho_{Ar},\rho_{can},N}^{(\theta)}(0)=\ln\det\square_{L,\rho_{Ar},h} -\ln{B_{g,k}}+\frac{1}{2}\ln \bigg(\frac{\det\Delta_{\rho_{Ar}}}{\vol(M,\rho_{Ar}) \det \mathfrak{Im}(\tau)}\bigg).
        \end{equation}
    \end{theorem}

\begin{theorem}\label{DifférenceLaplacienMagPotentiel}
    If $h$ is an admissible metric, $\hat{h}_{k,V}=e^{kV}h$ and $2\pi +\frac{1}{2}\Delta_{\rho_{can}}V>0$, then: 
    \begin{align*}
    \ln \det\square_{L,\rho_{Ar},\hat{h}_{k,V}}-\ln\det\square_{L,\rho_{Ar},h}&=k\left[\frac{1}{2}\int_{M}\mu_{\rho_{can}}\left(1+\frac{\Delta_{\rho_{can}}V}{4\pi} \right)\ln\left(1+\frac{\Delta_{\rho_{can}}V}{4\pi} \right)+\frac{1}{8\pi}\int_{M}\mu_{\rho_{can}}(R_{\rho_{can}}-8\pi(1-g))V \right]\\[1.5mm]&+\frac{2}{3}(1-g)\int_{M}\mu_{\rho_{can}}\ln\left(1+\frac{\Delta_{\rho_{can}}V}{4\pi} \right) +\frac{1}{24\pi}\int_{M}\mu_{\rho_{can}}\left(R_{\rho_{can}}-8\pi(1-g) \right)\ln \left(1+\frac{\Delta_{\rho_{can}}V}{4\pi} \right)\\[1.5mm]&+\frac{1}{48\pi}\int_{M}\mu_{\rho_{can}}\ln \left(1+\frac{\Delta_{\rho_{can}}V}{4\pi} \right)\Delta_{\rho_{can}}\ln\left(1+ \frac{\Delta_{\rho_{can}}V}{4\pi} \right)+O\left(\frac{1}{k} \right).
    \end{align*}
\end{theorem}

\begin{theorem}\label{DevAsympLapMagAdm}
    With $h$ an admissible metric 
    \begin{align*}
        \ln\det\square_{L,\rho_{Ar},h}&=-\frac{k}{2}\ln k +k\left(2\pi(1-g)S_{AY}(\hat{\phi}_{Ar},\rho_{can})-\frac{1}{4}S_{M}(-\sigma_{Ar},-\phi_{Ar},\rho_{Ar})-\frac{\ln2}{2}\right) +\frac{2(g-1)}{3}\ln k\\&- \frac{1-g}{12}(24\zeta'(-1)+2\ln2\pi+7+8\ln2)+\frac{\Psi_{P}(\rho_{can})-S_{L}(-\sigma_{Ar},\rho_{Ar})}{6\pi}-\frac{\Psi_{P}(\rho_{can})}{12\pi} +O\left(\frac{\ln k}{k} \right).
    \end{align*}
\end{theorem}

\begin{theorem}\label{TheoremeModifPreFinal}
    Let $L \to M$ be a line bundle of degree $k$, with $k=N+1-g$, $h$ be an admissible metric on $L$, and $V$ be a smooth potential such that: $2\pi+\frac{1}{2}\Delta_{\rho_{can}}V>0.$ Then, the modified partition function satisfies:
    \begin{align*}
    \ln\mathcal{Z}_{\rho_{Ar},{\rho_{can}},N}^{(\theta)}(V)=\ln \mathcal{Z}_{\rho_{Ar},{\rho_{can}},N}^{(\theta)}(0)+\ln \det\square_{L,\rho_{Ar},\hat{h}_{k,V}}-\ln \det\square_{L,\rho_{Ar},h}+k^{2}S_{\infty}(V)+\frac{k}{2}S_{1}(0,-V,\rho_{Ar},h)
    \end{align*} where $\ln \det\square_{L,\rho_{Ar},\hat{h}_{k,V}}-\ln \det\square_{L,\rho_{Ar},h}$ is given in Theorem \ref{DifférenceLaplacienMagPotentiel}, and
    \begin{align*}
        \ln \mathcal{Z}_{\rho_{Ar},\rho_{can},N}^{(\theta)}(0)&=-\frac{k}{2}\ln k +k\left(2\pi (1-g)S_{AY}(\hat{\phi}_{Ar},\rho_{can})-\frac{1}{4}S_{M}(-\sigma_{Ar},-\phi_{Ar},\rho_{Ar})+\ln2\pi -\frac{\ln2}{2}\right)+\frac{2(g-1)}{3}\ln k \\&-(1-g)\left(-4\zeta'(-1) -\frac{4}{3}\ln2\pi +\frac{5}{6}+\frac{\ln2}{6}\right)+\frac{\Psi_{P}(\rho_ {can})-S_{L}(-\sigma_{Ar},\rho_{Ar})}{6\pi}-\frac{\Psi_{P}(\rho_{can})}{12\pi}\\&+\frac{1}{2}\ln \left(\frac{\det\Delta_{\rho_{Ar}}}{\vol(M,\rho_{Ar})\det \mathfrak{Im}(\tau)} \right)+O\left(\frac{\ln k}{k} \right).
\end{align*}
\end{theorem}

\begin{proof}
    All we need to do to obtain our result is to combine Theorems \ref{TheoremPartitionGénéral} and \ref{DevAsympLapMagAdm}. 
    We have already proved the first equality (see Theorem \ref{TheoremPartitionGénéral}). We will prove the result for $\ln \mathcal{Z}_{\rho_{Ar},\rho_{can,N}}^{(\theta)}(0).$ We also know that (see Theorem \ref{TheoremPartitionGénéral}): $$\ln \mathcal{Z}_{\rho_{Ar},\rho_{can,N}}^{(\theta)}(0)=\ln\det\square_{L,\rho_{Ar},h}-\ln B_{g,k}+\frac{1}{2}\ln \left(\frac{\det \Delta_{\rho_{Ar}}}{\vol(M,\rho_{Ar})\det\mathfrak{Im}(\tau)} \right)$$ and (see Equation\eqref{lnBgk}) $$\ln B_{g,k}=-k\ln 2\pi +\frac{1-g}{2}\left(\frac{1}{2}-12\zeta'(-1)-3\ln2\pi-\ln2 \right)$$ and (see Theorem \ref{DevAsympLapMagAdm})   \begin{align*}
        \ln\det\square_{L,\rho_{Ar},h}&=-\frac{k}{2}\ln k +k\left(2\pi(1-g)S_{AY}(\hat{\phi}_{Ar},\rho_{can})-\frac{1}{4}S_{M}(-\sigma_{Ar},-\phi_{Ar},\rho_{Ar})-\frac{\ln2}{2}\right) +\frac{2(g-1)}{3}\ln k\\&- \frac{1-g}{12}(24\zeta'(-1)+2\ln2\pi+7+8\ln2)+\frac{\Psi_{P}(\rho_{can})-S_{L}(-\sigma_{Ar},\rho_{Ar})}{6\pi}-\frac{\Psi_{P}(\rho_{can})}{12\pi} +O\left(\frac{\ln k}{k} \right).
    \end{align*}
    Hence: 
    \begin{align*}
        \ln \mathcal{Z}_{\rho_{Ar},\rho_{can},N}^{(\theta)}(0)&=-\frac{k}{2}\ln k +k\left(2\pi (1-g)S_{AY}(\hat{\phi}_{Ar},\rho_{can})-\frac{1}{4}S_{M}(-\sigma_{Ar},-\phi_{Ar},\rho_{Ar})+\ln2\pi -\frac{\ln2}{2}\right)+\frac{2(g-1)}{3}\ln k \\&-(1-g)\left(-4\zeta'(-1) -\frac{4}{3}\ln2\pi +\frac{5}{6}+\frac{\ln2}{6}\right)+\frac{\Psi_{P}(\rho_ {can})-S_{L}(-\sigma_{Ar},\rho_{Ar})}{6\pi}-\frac{\Psi_{P}(\rho_{can})}{12\pi}\\&+\frac{1}{2}\ln \left(\frac{\det\Delta_{\rho_{Ar}}}{\vol(M,\rho_{Ar})\det \mathfrak{Im}(\tau)} \right)+O\left(\frac{\ln k}{k} \right).
    \end{align*}
\end{proof}

\end{document}
